\documentclass[11pt]{article}

% A synthetic seed paper for Lattice's research environment: plain article
% class, a hand-written bibliography, and code listings.
\usepackage[utf8]{inputenc}
\usepackage[T1]{fontenc}
\usepackage{lmodern}
\usepackage{amsmath,amssymb,amsthm}
\usepackage{listings}
\usepackage{booktabs}
\usepackage{url}

\newtheorem{theorem}{Theorem}
\newtheorem{definition}{Definition}
\newcommand{\Oh}{\mathcal{O}}

\lstset{basicstyle=\ttfamily\small, columns=fullflexible, frame=single}

\title{Notes on Reservoir Sampling}
\author{Seed Author\thanks{Lattice Research Environment, \texttt{seed@lattice.example}.}}
\date{}

\begin{document}
\maketitle

\begin{abstract}
These notes derive reservoir sampling, which keeps a uniform random sample of $k$ items from a stream of unknown length in one pass and $\Oh(k)$ memory.
We state the invariant that makes it correct, give a short implementation, and compare it with two alternatives.
\end{abstract}

\section{The problem}
\label{sec:problem}

A stream delivers items $x_1, x_2, \ldots$ one at a time, and its length $n$ is not known in advance.
We want, at every point, a sample of $k$ items such that every subset of size $k$ of the items seen so far is equally likely.
Storing the whole stream would solve the problem but needs $\Oh(n)$ memory; the algorithm below needs only the sample itself \cite{vitter1985random}.

\begin{definition}
\label{def:uniform}
A sample $R$ of size $k$ drawn from the first $t \ge k$ items is \emph{uniform} if $\Pr[x_i \in R] = k/t$ for every $i \le t$.
\end{definition}

\section{The algorithm}
\label{sec:algorithm}

Keep the first $k$ items.
When item $t > k$ arrives, draw $j$ uniformly from $\{1, \ldots, t\}$; if $j \le k$, replace the $j$-th item of the reservoir with $x_t$.
Listing~\ref{lst:reservoir} gives an implementation.

\begin{lstlisting}[language=Python, caption={Reservoir sampling in Python.}, label={lst:reservoir}]
import random

def reservoir(stream, k):
    sample = []
    for t, item in enumerate(stream, start=1):
        if t <= k:
            sample.append(item)
        else:
            j = random.randint(1, t)
            if j <= k:
                sample[j - 1] = item
    return sample
\end{lstlisting}

\begin{theorem}
\label{thm:uniform}
After $t \ge k$ items, the reservoir is a uniform sample in the sense of Definition~\ref{def:uniform}.
\end{theorem}

\begin{proof}
By induction on $t$. For $t = k$ every item is kept, so $\Pr[x_i \in R] = 1 = k/t$.
Suppose the claim holds after $t - 1$ items. Item $x_t$ enters with probability $k/t$.
An earlier item stays if it was in the reservoir and is not the one replaced, which happens with probability
\[
  \frac{k}{t-1}\left(1 - \frac{k}{t}\cdot\frac{1}{k}\right) = \frac{k}{t-1}\cdot\frac{t-1}{t} = \frac{k}{t}. \qedhere
\]
\end{proof}

\section{Alternatives}
\label{sec:alternatives}

Table~\ref{tab:alternatives} compares the algorithm with two other ways to sample.
Theorem~\ref{thm:uniform} is what the first row relies on.

\begin{table}[h]
  \centering
  \caption{Three ways to keep a uniform sample of $k$ items from $n$.}
  \label{tab:alternatives}
  \begin{tabular}{@{}lll@{}}
    \toprule
    Method & Memory & Needs $n$ in advance \\
    \midrule
    Reservoir sampling & $\Oh(k)$ & no \\
    Store everything, then sample & $\Oh(n)$ & no \\
    Selection sampling \cite{knuth1997art} & $\Oh(k)$ & yes \\
    \bottomrule
  \end{tabular}
\end{table}

Two refinements are worth knowing:
\begin{enumerate}
  \item skipping ahead by a random number of items, so random numbers are drawn only when an item enters \cite{vitter1985random};
  \item weighted sampling, where item $i$ enters with probability proportional to a weight $w_i$ \cite{efraimidis2006weighted}.
\end{enumerate}
The first is item~1 above and keeps the guarantee of Section~\ref{sec:algorithm}.

\begin{thebibliography}{9}
\bibitem{vitter1985random}
J.~S. Vitter.
\newblock Random sampling with a reservoir.
\newblock \emph{ACM Transactions on Mathematical Software}, 11(1):37--57, 1985.

\bibitem{knuth1997art}
D.~E. Knuth.
\newblock \emph{The Art of Computer Programming, Volume 2: Seminumerical Algorithms}.
\newblock Addison-Wesley, 3rd edition, 1997.

\bibitem{efraimidis2006weighted}
P.~S. Efraimidis and P.~G. Spirakis.
\newblock Weighted random sampling with a reservoir.
\newblock \emph{Information Processing Letters}, 97(5):181--185, 2006.
\end{thebibliography}

\end{document}
